\documentclass[conference]{IEEEtran}
\IEEEoverridecommandlockouts

\usepackage{cite}
\usepackage{amsmath,amssymb,amsfonts}
\usepackage{graphicx}
\usepackage{textcomp}
\usepackage{xcolor}
\usepackage{mathtools}
\usepackage{booktabs}

\def\BibTeX{{\rm B\kern-.05em{\sc i\kern-.025em b}\kern-.08em
T\kern-.1667em\lower.7ex\hbox{E}\kern-.125emX}}

\newcommand{\E}{\mathbb{E}}

\begin{document}

\title{Outage Performance and Deep Reinforcement Learning-Based Orientation
Optimization of OIRS-Aided VLC Systems under Realistic UE Orientation
Constraints}

\maketitle

\begin{abstract}
Optical intelligent reflecting surface (OIRS)-aided visible light
communication (VLC) systems effectively mitigate complete line-of-sight (LoS)
blockages, yet their reliability remains compromised by random user equipment
(UE) orientations. This paper develops a rigorous analytical and learning
framework for such links. Adopting a measurement-based truncated-Laplace
orientation model, we derive the exact moments of the incidence cosine and
approximate the equivalent non-line-of-sight (NLoS) channel gain with a Gamma
distribution via exact moment matching. Closed-form outage probability
expressions follow from the Meijer G-function, and we prove a strict
``zero-outage'' signal-to-noise ratio (SNR) cliff dictated by the physical
field-of-view bound---a hard transition absent under conventional uniform
orientation assumptions. Since the OIRS yaw and roll jointly serving $K$ users
render the expected achievable data rate (ADR) maximization non-convex with a
stochastic objective, we cast the orientation design as a continuous-action
reinforcement learning problem and comparatively evaluate five deep
reinforcement learning (DRL) algorithms---DDPG, SAC, PPO, CAPG, and the
proposed ECDPO---under identical conditions. Off-policy agents recover the
expected-ADR optimum to within $0.5$ of an exhaustive grid reference, and
retraining with a min--max Age of Information (AoI) reward restores fairness
to users abandoned by throughput-optimal steering, attaining the theoretical
minimum AoI for all users. Analytical results are validated by Monte Carlo
simulations throughout.
\end{abstract}

\begin{IEEEkeywords}
Deep reinforcement learning, Meijer G-function, optical intelligent
reflecting surface, outage probability, random UE orientation.
\end{IEEEkeywords}

\section{Introduction}
Visible light communication (VLC) is a cornerstone technology for indoor
Internet of Things (IoT) networks, offering massive unlicensed bandwidth and
spatial security. Its performance, however, is bottlenecked by line-of-sight
(LoS) blockages and the stochastic geometry of random user equipment (UE)
orientations \cite{eroglu2018}. Optical intelligent reflecting surfaces
(OIRSs) bypass blockages by establishing robust non-line-of-sight (NLoS)
paths \cite{aboagye2023}, but two gaps persist in their analysis. First,
existing frameworks frequently assume uniformly distributed UE orientation,
which assigns equal probability to physically improbable incidence angles,
artificially inflates the channel variance, and pulls the minimum channel
gain toward zero, yielding overly pessimistic reliability predictions.
Second, the orientation of the OIRS itself---a common yaw $\gamma$ and roll
$\omega$ serving all users---must be \emph{optimized} against an objective
that is non-convex and, under random UE orientation, observable only through
noisy samples, precluding closed-form solutions.

This paper addresses both gaps. Our contributions are:
\begin{itemize}
\item We derive the exact centralized moments of the incidence cosine under a
measurement-based truncated-Laplace model \cite{soltani2018} and approximate
the multi-element NLoS gain with a Gamma distribution via exact moment
matching, strictly enforcing the receiver field-of-view (FoV) bounds.
\item We obtain exact closed-form outage probability expressions using the
Meijer G-function with intensity-modulation/direct-detection (IM/DD) capacity
bounds, and prove a ``zero-outage cliff'': an SNR beyond which the outage
probability is exactly zero, in contrast to the polynomial tails of
unbounded fading models.
\item We formulate the OIRS orientation design as a stochastic optimization
of the \emph{expected} achievable data rate (ADR) and solve it with five
continuous-action deep reinforcement learning (DRL) algorithms---DDPG
\cite{lillicrap2016}, SAC \cite{haarnoja2018}, PPO \cite{schulman2017}, CAPG,
and the proposed enhanced continuous deep policy optimization
(ECDPO)---trained under identical conditions and benchmarked against a
common-random-number grid search; a min--max Age of Information (AoI) reward
variant demonstrates cross-layer fairness control.
\end{itemize}

\section{System Model}
We consider an indoor VLC system with an $M$-LED access point and $K$
photodetector (PD)-equipped UEs. The direct LoS links are completely
obstructed; communication is supported entirely by a wall-mounted OIRS of $N$
elements acting as a unified reflector with common yaw $\gamma$ and roll
$\omega$. The equivalent NLoS channel gain of user $k$ is
\begin{equation}
Z_k = a_k \sum_{l=1}^{M}\sum_{i=1}^{N} \cos\bigl(\psi_k^{i}\bigr),
\label{eq:Zk}
\end{equation}
where $\psi_k^{i}$ is the random incidence angle of the path through the
$i$th element and $a_k$ aggregates the deterministic path loss, Lambertian
emission, and receiver gains. The OIRS orientation enters through the
unified-reflector steering: with mirror normal
$\mathbf n(\gamma,\omega)=\mathbf R_y(\omega)\mathbf R_z(\gamma)(1,0,0)^{\mathsf T}$
and specular direction $\mathbf d_{\mathrm{ref}}$,
\begin{equation}
a_k(\gamma,\omega)=\rho\,\frac{(m+1)A_{\mathrm{pd}}}{2\pi (d_1+d_{2,k})^2}
\cos^m\!\phi\,
\bigl(\mathbf d_{\mathrm{ref}}^{\mathsf T}\mathbf u_k\bigr)_{+}^{\kappa},
\label{eq:ak}
\end{equation}
where $d_1$ ($d_{2,k}$) is the LED--OIRS (OIRS--user) distance, $\phi$ the
LED radiance angle, $\mathbf u_k$ the OIRS$\to$user direction, and $\kappa$
the reflector beam-concentration exponent.

\subsection{Truncated-Laplace Orientation Statistics}
Define $\xi=\cos(\psi_k^{i})$. Constraining the orientation to practical
human usage, $\xi$ follows the truncated-Laplace law
\begin{equation}
f_{\xi}(\xi) = \frac{1}{\Lambda}\exp\!\left(-\frac{|\xi-\mu_\xi|}{b_\xi}\right),
\quad \xi\in[\xi_{\min},\xi_{\max}],
\label{eq:laplace}
\end{equation}
where $\mu_\xi$, $b_\xi$ follow from the measured polar-angle statistics
\cite{soltani2018} and $\xi_{\min}$ is dictated by the receiver FoV. With
$u=\xi-\mu_\xi$, $\alpha=\xi_{\min}-\mu_\xi\le 0$,
$\beta=\xi_{\max}-\mu_\xi\ge 0$, splitting the integrals at $u=0$ gives
\begin{equation}
\Lambda = b_\xi\bigl[2 - e^{\alpha/b_\xi} - e^{-\beta/b_\xi}\bigr],
\end{equation}
and, by (repeated) integration by parts, the exact moments
\begin{equation}
\E[u] = \frac{b_\xi}{\Lambda}\Bigl[e^{\alpha/b_\xi}(b_\xi-\alpha)
-e^{-\beta/b_\xi}(\beta+b_\xi)\Bigr],
\label{eq:moment1}
\end{equation}
\begin{equation}
\begin{aligned}
\E[u^2] &= \frac{1}{\Lambda}\Bigl[\,4b_\xi^3
- e^{\alpha/b_\xi}\bigl(b_\xi\alpha^2-2b_\xi^2\alpha+2b_\xi^3\bigr)\\
&\quad\; - e^{-\beta/b_\xi}\bigl(b_\xi\beta^2+2b_\xi^2\beta+2b_\xi^3\bigr)
\Bigr],
\end{aligned}
\label{eq:moment2}
\end{equation}
whence $\E[\xi]=\mu_\xi+\E[u]$ and
$\mathrm{Var}[\xi]=\E[u^2]-\E[u]^2$.

\subsection{Multi-Element Gamma Approximation}
The exact PDF of $Z_k$ requires an $MN$-fold convolution of
\eqref{eq:laplace}, which is intractable in closed form. Matching the exact
mean $\E[Z_k]=MNa_k\E[\xi]$ and variance
$\mathrm{Var}[Z_k]=MNa_k^2\mathrm{Var}[\xi]$ to a Gamma law
\begin{equation}
f_{Z_{k}}(z) = \frac{z^{k_{g_k}-1}e^{-z/\theta_{g_k}}}
{\theta_{g_k}^{k_{g_k}}\Gamma(k_{g_k})},
\quad z \in [L_B,\, MNa_k\xi_{\max}],
\label{eq:gamma}
\end{equation}
with $L_B=MNa_k\xi_{\min}$ the strict physical lower bound, yields
\begin{equation}
k_{g_k} = MN\frac{\E[\xi]^2}{\mathrm{Var}[\xi]}, \qquad
\theta_{g_k} = a_k\frac{\mathrm{Var}[\xi]}{\E[\xi]}.
\label{eq:matching}
\end{equation}

\section{Outage Analysis and the Zero-Outage Cliff}
For IM/DD OOK signaling with target rate $R$, the instantaneous electrical
SNR is $\gamma_k=|Z_k|^2\bar{\gamma}$ and an outage occurs when $Z_k$ falls
below
\begin{equation}
z_{th} = \sqrt{\frac{2\pi}{e\,\bar{\gamma}} \bigl( 2^{2R} - 1 \bigr)}.
\label{eq:zth}
\end{equation}
Evaluating the Gamma CDF strictly from the physical bound $L_B$ and using the
Meijer G-representation of the incomplete gamma function (avoiding
infinite-series truncation errors),
\begin{equation}
\begin{aligned}
P_{out} &= \frac{1}{\Gamma(k_{g_k})} \Biggl[
G_{1, 2}^{1, 1}\!\left( \frac{z_{th}}{\theta_{g_k}} \,\middle|\,
\begin{matrix} 1 \\ k_{g_k},\, 0 \end{matrix} \right)\\[2pt]
&\qquad\qquad
- G_{1, 2}^{1, 1}\!\left( \frac{L_B}{\theta_{g_k}} \,\middle|\,
\begin{matrix} 1 \\ k_{g_k},\, 0 \end{matrix} \right) \Biggr].
\end{aligned}
\label{eq:pout}
\end{equation}
\emph{Zero-outage cliff:} unlike Rayleigh/Rician fading, whose support
extends to zero and whose outage decays polynomially, the truncated-Laplace
constraint enforces $L_B>0$. Since $z_{th}$ in \eqref{eq:zth} decreases
monotonically in $\bar{\gamma}$, the condition $z_{th}\le L_B$ is met at
finite power, whereupon the integral in \eqref{eq:pout} spans a null set and
$P_{out}$ equals \emph{exactly} zero. Equating $z_{th}=L_B$,
\begin{equation}
\bar{\gamma}_{\mathrm{cliff}} = \frac{2\pi}{e}\,
\frac{2^{2R} - 1}{\left( M N a_k \xi_{\min} \right)^2}.
\label{eq:cliff}
\end{equation}
Under the conventional uniform assumption with support pulled to
$\xi_{\min}\!\to\!0$, $\bar{\gamma}_{\mathrm{cliff}}\!\to\!\infty$ and no such
transition exists.

\section{DRL-Based OIRS Orientation Optimization}
The orientation design problem is
\begin{equation}
(\mathrm{P1}):\;\;
\max_{\gamma,\,\omega}\;
\E_{\xi}\!\left[\frac{1}{K}\sum_{k=1}^{K}
\frac{B}{2}\log_{2}\!\left(1+\frac{e}{2\pi}
\Bigl(\tfrac{P R_{\mathrm{PD}} Z_k}{Q}\Bigr)^{2}\right)\right]
\label{eq:p1}
\end{equation}
subject to $\gamma,\omega\in[-\pi/2,\pi/2]$, where the expectation is over
the random UE orientations in \eqref{eq:Zk}. Because \eqref{eq:p1} is
non-convex in $(\gamma,\omega)$ through \eqref{eq:ak} and only noisy samples
of the objective are observable, we solve it by continuous-action DRL:
single-step episodes in which the agent observes the (fixed, normalized)
geometry, outputs $\mathbf a=[\gamma,\omega]$, and receives one fresh draw of
all $MN$ incidence cosines per user, i.e., the reward is the sampled mean ADR.
Critics therefore learn the \emph{expected}-ADR surface.

Five algorithms are compared under identical state, action space, reward,
network width (two hidden layers of 64), seed, and budget (6{,}000
interactions): \textbf{DDPG} \cite{lillicrap2016} (deterministic actor,
Gaussian exploration $\sigma=0.3$); \textbf{SAC} \cite{haarnoja2018}
($\tanh$-Gaussian actor, twin soft critics, and \emph{automatic} entropy
temperature with target $-\dim\mathcal{A}$, which proved essential under the
noisy reward); \textbf{PPO} \cite{schulman2017} (clipped surrogate, on-policy
batches of 256 with KL-triggered early stopping and decaying entropy bonus);
\textbf{CAPG} (continuous-action likelihood-ratio policy gradient with
moving-average baseline and floored log-std); and \textbf{ECDPO}, an enhanced
continuous deep policy optimizer employing twin critics with a minimum
operator against noise-driven overestimation, delayed actor updates, Huber
critic loss, and exponentially decaying exploration noise
($0.5\!\to\!0.05$) for coarse-to-fine alignment. A reward-mode switch
replaces the ADR objective with the cross-layer min--max HARQ-IR AoI
objective $\min_{\gamma,\omega}\max_k \bar{\Delta}_{k,\mathrm{IR}}$,
computed in closed form via \eqref{eq:matching} and the Gamma CDF.

\section{Numerical Results}
A $5\times5\times3$~m room is simulated with $M=1$, $N=4$, $K=3$
users, $b_\xi=0.036$, $\xi_{\max}=0.9574$ \cite{soltani2018},
$B=20$~MHz, $P=3$~W, $R_{\mathrm{PD}}=0.4$~A/W, $\rho=0.95$, and
$\kappa=20$ for the throughput study. UE-orientation randomness is drawn by
accept--reject sampling of \eqref{eq:laplace}; all analytical curves are
validated by Monte Carlo (MC) simulation.

\begin{table}[t]
\centering
\caption{Learned orientations after 6{,}000 stochastic interactions.}
\label{tab:results}
\begin{tabular}{lcccc}
\toprule
Algorithm & $\gamma$ ($^\circ$) & $\omega$ ($^\circ$) &
$\E[\mathrm{ADR}]$ (Mbit/s) & \% of ref. \\
\midrule
DDPG & $16.54$ & $-7.82$ & $36.94$ & $99.5$ \\
SAC & $16.87$ & $-8.83$ & $37.11$ & $100.0$ \\
PPO & $6.78$ & $-8.72$ & $22.33$ & $60.2$ \\
CAPG & $8.30$ & $-7.48$ & $25.10$ & $67.6$ \\
ECDPO & $16.92$ & $-8.37$ & $37.06$ & $99.8$ \\
\midrule
Grid reference & $16.54$ & $-8.77$ & $37.12$ & $100$ \\
\bottomrule
\end{tabular}
\end{table}

\begin{figure}[t]
\centering
\includegraphics[width=0.98\columnwidth]{fig1_convergence.pdf}
\caption{DRL convergence on the stochastic OIRS-VLC channel (moving average
of the sampled mean ADR) versus the common-random-number grid reference.}
\label{fig:conv}
\end{figure}

\emph{Orientation optimization:} Table~\ref{tab:results} and
Fig.~\ref{fig:conv} show that the off-policy agents (SAC, ECDPO, DDPG)
recover the expected-ADR optimum to within $0.5$, with SAC's automatic
temperature yielding the closest match. The on-policy agents (PPO, CAPG)
settle on a shoulder of the multi-user landscape at this budget---a
sample-efficiency gap consistent with their discard-after-update use of the
noisy reward, not an algorithmic defect.

\begin{figure}[t]
\centering
\includegraphics[width=0.98\columnwidth]{fig4_outage.pdf}
\caption{Outage probability vs.\ transmit SNR at the ECDPO-learned
orientation ($R=2$, $\xi_{\min}=0.2$): analytical \eqref{eq:pout} vs.\
MC markers. The truncated-Laplace model reaches any target outage at far
lower SNR than the bounded uniform model and both hit the exact zero-outage
cliff \eqref{eq:cliff} at $146$~dB, whereas the conventional
$\xi_{\min}\!\to\!0$ uniform assumption exhibits no cliff.}
\label{fig:outage}
\end{figure}

\emph{Outage validation and cliff:} Fig.~\ref{fig:outage} confirms tight
agreement between \eqref{eq:pout} and MC throughout the waterfall region. At
$\bar{\gamma}_{\mathrm{cliff}}\!\approx\!146$~dB, given by \eqref{eq:cliff},
the analytical and empirical outage drop \emph{exactly} to zero; the
truncated-Laplace statistics reach $P_{out}{=}10^{-6}$ roughly $8$~dB before
the bounded uniform model, quantifying the pessimism of uniform bounds, while
the $\xi_{\min}\!\to\!0$ uniform tail never terminates.

\begin{figure}[t]
\centering
\includegraphics[width=0.98\columnwidth]{fig3_fairness.pdf}
\caption{Per-user expected ADR (left) and average HARQ-IR AoI (right) at the
ADR-optimal vs.\ AoI-optimal learned orientations.}
\label{fig:fair}
\end{figure}

\emph{Cross-layer fairness:} the throughput-optimal orientation abandons
UE1 entirely ($\E[\mathrm{ADR}_1]=0$, worst-case AoI $4.5$ rounds),
evidencing the unfairness of pure rate maximization. Retraining ECDPO with
the min--max AoI reward on a wide-beam configuration ($\kappa=6$,
$\bar{\gamma}=3\times10^{15}$) attains per-user AoI
$[1.5,\,1.5,\,1.5]$---the theoretical minimum, with every user decoding in
the first HARQ round---matching the min--max grid reference exactly
(Fig.~\ref{fig:fair}).

\section{Conclusion}
This paper unified a rigorous outage analysis of OIRS-aided VLC under
realistic truncated-Laplace UE orientation with a DRL-based optimization of
the OIRS orientation itself. Exact moment matching yielded a tractable Gamma
channel characterization and finite Meijer G-function outage expressions, and
the physical FoV bound was proven to induce an exact zero-outage SNR cliff
unreachable under conventional uniform assumptions. Framing the orientation
design as stochastic expected-ADR maximization, off-policy DRL agents
recovered the grid-search optimum to within $0.5$, and a min--max AoI
reward variant restored fairness to misaligned users at the theoretical AoI
minimum. These results establish OIRS-aided VLC links as substantially more
deterministic and controllable than uniform bounds suggest, providing a
foundation for URLLC-grade MAC design.

\bibliographystyle{IEEEtran}
\begin{thebibliography}{00}

\bibitem{aboagye2023}
S. Aboagye \emph{et al.}, ``RIS-assisted visible light communication systems:
A tutorial,'' \emph{IEEE Commun. Surv. Tutor.}, vol. 25, no. 1, pp. 251--288,
2023.

\bibitem{eroglu2018}
Y. S. Ero\u{g}lu \emph{et al.}, ``Impact of random receiver orientation on
visible light communications channel,'' \emph{IEEE Trans. Commun.}, vol. 67,
no. 2, pp. 1313--1325, 2019.

\bibitem{soltani2018}
M. D. Soltani \emph{et al.}, ``Modeling the random orientation of mobile
devices,'' \emph{IEEE Trans. Commun.}, vol. 67, no. 3, pp. 2157--2172, 2019.

\bibitem{lillicrap2016}
T. P. Lillicrap \emph{et al.}, ``Continuous control with deep reinforcement
learning,'' in \emph{Proc. ICLR}, 2016.

\bibitem{haarnoja2018}
T. Haarnoja \emph{et al.}, ``Soft actor-critic: Off-policy maximum entropy
deep reinforcement learning with a stochastic actor,'' in \emph{Proc. ICML},
2018, pp. 1861--1870.

\bibitem{schulman2017}
J. Schulman \emph{et al.}, ``Proximal policy optimization algorithms,''
\emph{arXiv:1707.06347}, 2017.

\end{thebibliography}

\end{document}
